\documentclass[a4paper,fontsize = 7pt]{scrartcl}
\usepackage[utf8]{inputenc}
\usepackage[german]{babel}

\usepackage{lipsum}

\usepackage{tikz}

\usetikzlibrary{positioning, calc, fit, matrix, decorations, angles, quotes, babel, 3d, arrows.meta, shadings, shapes.geometric}

\tikzset{>=stealth,
	overlaynote/.style={font = \footnotesize, fill = overlaycolor, draw = overlaycolor, rounded corners = 3pt},
	overlayarrow/.style={thick, gray, ->}
}

% 3 column landscape layout with fewer margins
\usepackage[landscape, left=0.75cm, top=1cm, right=0.75cm, bottom=1.5cm, footskip=5pt]{geometry}
\usepackage{flowfram}
\ffvadjustfalse
\setlength{\columnsep}{1cm}
\Ncolumn{4}

% define nice looking boxes
\usepackage[most]{tcolorbox}

% a base set, that is then customised
\tcbset {
  base/.style={
    boxrule=0.3mm,
    leftrule=1mm,
    left=1.75mm,
    arc=0mm,
    fonttitle=\bfseries,
    colbacktitle=black!10!white,
    coltitle=black,
    toptitle=0.75mm,
    bottomtitle=0.25mm,
    title={#1}
  }
}

\definecolor{nicegreen}{rgb}{0.004, 0.5, 0.35}
\newtcolorbox{mainbox}[1]{
  colframe=nicegreen,
  base={#1}
}

\newtcolorbox{subbox}[1]{
  colframe=black!20!white,
  base={#1}
}

% Mathematical typesetting & symbols
\usepackage{amsthm, mathtools, amssymb, centernot}
\usepackage{marvosym, wasysym}
\allowdisplaybreaks

% Tables
\usepackage{tabularx, multirow}
\usepackage{booktabs}
\renewcommand*{\arraystretch}{2}

% Make enumerations more compact
\usepackage{enumitem}
\setitemize{itemsep=0.5pt}
\setenumerate{itemsep=0.75pt}

% To include sketches & PDFs
\usepackage{graphicx}


% For hyperlinks
\usepackage{hyperref}
\usepackage{amsopn}
\usepackage{color}
\usepackage{amsfonts}
\usepackage{amsmath}
\usepackage{amstext}
\hypersetup{
  colorlinks=true
}

% Metadata
\title{Cheatsheet ANA-G}
\author{Thomas Schnyder}
\date{Januar 2026}

% Math helper stuff
\def\limn{\lim_{n\to \infty}}
\def\limxo{\lim_{x\to 0}}
\def\limxi{\lim_{x\to\infty}}
\def\limxn{\lim_{x\to-\infty}}
\def\sumk{\sum_{k=1}^\infty}
\def\sumn{\sum_{n=0}^\infty}
\def\Z{\mathbb{Z}}
\def\R{\mathbb{R}}
\def\N{\mathbb{N}}
\def\C{\mathbb{C}}
\def\dx{\text{ d}x}
\def\oof{\mathcal{O}}

%removing paragraph indent
\setlength{\parindent}{0cm}
\setcounter{tocdepth}{4}
\setcounter{secnumdepth}{4}

\begin{document}

\maketitle

\section{Differenzen- und Differentialquotient}

\begin{subbox}{Steigung einer Geraden}
	\textbf{Gegeben:} Zwei Punkte $P_1(x_1/y_1)$ und $P_2(x_2/y_2)$ \\
	\textbf{Steigung:}
	$$m = \frac{\Delta y}{\Delta x} = \frac{y_2 - y_1}{x_2 - x_1} = \frac{y_1 - y_2}{x_1 - x_2}$$
\end{subbox}

\begin{mainbox}{Geradengleichungen}
\textbf{Allgemeine Form:} Gerade mit Steigung $m$ und $y$-Achsenabschnitt $t$
$$y(x) = m \cdot x + t$$

\textbf{Punkt-Steigungsform:} Gerade durch $P_0(x_0/y_0)$ mit Steigung $m$
$$y(x) = m \cdot (x - x_0) + y_0 = m \cdot x + (y_0 - m \cdot x_0)$$

\textbf{Zwei-Punkte-Form:} Gerade durch $P_1(x_1/y_1)$ und $P_2(x_2/y_2)$
	$$y(x) = \frac{y_2 - y_1}{x_2 - x_1} \cdot (x - x_1) + y_1$$
\end{mainbox}

\subsection{Differenzenquotient (Sekantensteigung)}

\begin{mainbox}{Sekantensteigung - Form 1}
	\textbf{Gegeben:} Funktion $f(x)$, zwei Stellen $x_1$ und $x_2$ \\
	\textbf{Gesucht:} Steigung der Sekante durch $P_1(x_1/f(x_1))$ und $P_2(x_2/f(x_2))$ \\
	\textbf{Differenzenquotient:}
	$$m = \frac{\Delta y}{\Delta x} = \frac{f(x_2) - f(x_1)}{x_2 - x_1} = \frac{f(x_1) - f(x_2)}{x_1 - x_2}$$
\end{mainbox}

\begin{mainbox}{Sekantensteigung - Form 2}
	\textbf{Gegeben:} Funktion $f(x)$, Stelle $x_0$ und Abstand $h$ \\
	\textbf{Gesucht:} Steigung der Sekante durch $P_0(x_0/f(x_0))$ und $P(x_0+h/f(x_0+h))$ \\
	\textbf{Differenzenquotient:}
	$$m = \frac{\Delta y}{\Delta x} = \frac{f(x_0 + h) - f(x_0)}{h}$$
	\textbf{Interpretation:} Durchschnittliche Änderungsrate der Funktion zwischen $x_0$ und $x_0 + h$
\end{mainbox}

\begin{subbox}{Beispiel: Sekantensteigung}
	Für $f(x) = x^2$ mit $x_1 = -1$ und $x_2 = 3$:
	$$m = \frac{f(3) - f(-1)}{3 - (-1)} = \frac{9 - 1}{4} = \frac{8}{4} = 2$$
\end{subbox}

\subsection{Differentialquotient (Tangentensteigung)}

\begin{mainbox}{Ableitung einer Funktion an der Stelle $x_0$ - Form 1}
	\textbf{Gegeben:} Funktion $f(x)$, feste Stelle $x_0$ \\
	\textbf{Gesucht:} Steigung der Tangente an $f$ bei $x_0$ \\
	\textbf{Differentialquotient (Ableitung):}
	$$f'(x_0) = \lim_{x \to x_0} \frac{f(x) - f(x_0)}{x - x_0}$$
	\textbf{Interpretation:} Momentane Änderungsrate der Funktion an der Stelle $x_0$
\end{mainbox}

\begin{mainbox}{Ableitung einer Funktion an der Stelle $x_0$ - Form 2}
	\textbf{Gegeben:} Funktion $f(x)$, feste Stelle $x_0$ \\
	\textbf{Differentialquotient (Ableitung):}
	$$f'(x_0) = \lim_{h \to 0} \frac{f(x_0 + h) - f(x_0)}{h}$$
	\textbf{Zusammenhang:} Grenzwert der Sekantensteigung für $h \to 0$
\end{mainbox}

\begin{subbox}{Beispiel: Ableitung von $f(x) = x^2$ bei $x_0 = 2$}
	\begin{align*}
		f'(2) &= \lim_{h \to 0} \frac{(2+h)^2 - 2^2}{h} \\
		&= \lim_{h \to 0} \frac{4 + 4h + h^2 - 4}{h} \\
		&= \lim_{h \to 0} \frac{4h + h^2}{h} \\
		&= \lim_{h \to 0} (4 + h) = 4
	\end{align*}
\end{subbox}

\subsection{Ableitungsfunktion}

\begin{mainbox}{Ableitungsfunktion $f'(x)$}
	Die Ableitungsfunktion $f'(x)$ ordnet jeder Stelle $x$ die Tangentensteigung an dieser Stelle zu:
	$$f'(x) = \lim_{h \to 0} \frac{f(x + h) - f(x)}{h}$$
	\textbf{Alternative Schreibweisen:}
	$$\frac{df}{dx}, \quad \frac{d}{dx}f(x), \quad Df(x)$$
\end{mainbox}

\begin{mainbox}{Wichtige Ableitungsregeln}
\textbf{Potenzregel:}
$$f(x) = x^n \quad \Rightarrow \quad f'(x) = n \cdot x^{n-1}$$

\textbf{Spezialfälle:}
\begin{itemize}
	\item $f(x) = x^3 \Rightarrow f'(x) = 3x^2$
	\item $f(x) = x^5 \Rightarrow f'(x) = 5x^4$
	\item $f(x) = \frac{1}{x} = x^{-1} \Rightarrow f'(x) = -x^{-2} = -\frac{1}{x^2}$
\end{itemize}

\textbf{Trigonometrische Funktionen:}
\begin{itemize}
	\item $f(x) = \sin(x) \Rightarrow f'(x) = \cos(x)$
	\item $f(x) = \cos(x) \Rightarrow f'(x) = -\sin(x)$
\end{itemize}

\textbf{Konstante:}
$$f(x) = c \quad \Rightarrow \quad f'(x) = 0$$
\end{mainbox}

\subsection{Ableitungsregeln}

\begin{mainbox}{Linearität der Ableitung}
Die Ableitung ist linear, d.h.\ sie ist verträglich mit Addition/Subtraktion sowie Multiplikation mit Konstanten.

\textbf{Summen-/Differenzregel:}
$$(f(x) \pm g(x))' = f'(x) \pm g'(x)$$
Die Ableitung einer Summe (Differenz) ist die Summe (Differenz) der Ableitungen.

\textbf{Konstanter Faktor:}
$$(c \cdot f(x))' = c \cdot f'(x)$$
Einen konstanten Faktor darf man vor die Ableitung ziehen.

\textbf{Konstanter Summand:}
$$(f(x) + c)' = f'(x)$$
Ein konstanter Summand entfällt beim Ableiten.
\end{mainbox}

\begin{subbox}{Beispiel: Polynom ableiten}
	$f(x) = 3x^2 - 5x + 7$
	\begin{align*}
		f'(x) &= (3x^2 - 5x + 7)' \\
		&= (3x^2)' - (5x)' + (7)' \\
		&= 3 \cdot 2x - 5 \cdot 1 + 0 \\
		&= 6x - 5
	\end{align*}
\end{subbox}

\begin{mainbox}{Produktregel}
$$(f(x) \cdot g(x))' = f'(x) \cdot g(x) + f(x) \cdot g'(x)$$
Die Ableitung wandert von einem Faktor zum anderen.

\textbf{Für drei Faktoren:}
$$(f(x) \cdot g(x) \cdot h(x))' = f'(x) \cdot g(x) \cdot h(x) + f(x) \cdot g'(x) \cdot h(x) + f(x) \cdot g(x) \cdot h'(x)$$

\textbf{Allgemein für $n$ Faktoren:}
$$(p \cdot q \cdot r \cdot s)' = p' \cdot q \cdot r \cdot s + p \cdot q' \cdot r \cdot s + p \cdot q \cdot r' \cdot s + p \cdot q \cdot r \cdot s'$$
\end{mainbox}

\begin{subbox}{Beispiel: Produktregel}
	$f(x) = x \cdot e^x$
	\begin{align*}
		f'(x) &= (x \cdot e^x)' \\
		&= (x)' \cdot e^x + x \cdot (e^x)' \\
		&= 1 \cdot e^x + x \cdot e^x \\
		&= (1 + x) \cdot e^x
	\end{align*}
\end{subbox}

\begin{mainbox}{Quotientenregel}
$$\left(\frac{f(x)}{g(x)}\right)' = \frac{f'(x) \cdot g(x) - f(x) \cdot g'(x)}{g^2(x)}$$
\textbf{Merkregel:} NAZ $-$ ZAN durch N$^2$ (Nenner mal Ableitung Zähler minus Zähler mal Ableitung Nenner, durch Nenner zum Quadrat)
\end{mainbox}

\begin{subbox}{Beispiel: Quotientenregel}
	$f(x) = \frac{x}{e^x}$
	\begin{align*}
		f'(x) &= \left(\frac{x}{e^x}\right)' = \frac{(x)' \cdot e^x - x \cdot (e^x)'}{(e^x)^2} \\
		&= \frac{1 \cdot e^x - x \cdot e^x}{e^{2x}} = \frac{(1-x) \cdot e^x}{e^{2x}} = \frac{1-x}{e^x}
	\end{align*}
\end{subbox}

\begin{mainbox}{Kettenregel}
$$(f(g(x)))' = f'(g(x)) \cdot g'(x)$$
Eine zusammengesetzte Funktion wird differenziert, indem man die äussere Funktion ableitet und mit der Ableitung der inneren Funktion multipliziert.

\textbf{Merkregel:} Äussere Ableitung mal innere Ableitung (``Innere Funktion nachdifferenzieren'')

\textbf{Leibniz-Notation:}
$$\frac{d(f \circ g)}{dx} = \frac{df}{dg} \cdot \frac{dg}{dx}$$
\end{mainbox}

\begin{subbox}{Beispiel: Kettenregel}
	$f(x) = \ln(x^2 + 1)$

	Setze $g(y) = \ln(y)$ mit $g'(y) = \frac{1}{y}$ und $h(x) = x^2 + 1$ mit $h'(x) = 2x$.
	\begin{align*}
		f'(x) &= g'(h(x)) \cdot h'(x) \\
		&= \frac{1}{x^2 + 1} \cdot 2x = \frac{2x}{x^2 + 1}
	\end{align*}
\end{subbox}

\begin{mainbox}{Ableitung der Umkehrfunktion}
Für eine umkehrbare Funktion $f$ mit Umkehrfunktion $f^{-1}$ gilt:
$$(f^{-1}(x))' = \frac{1}{f'(f^{-1}(x))}$$

\textbf{Herleitung:} Aus $f(f^{-1}(x)) = x$ folgt durch Ableiten mit Kettenregel:
$$f'(f^{-1}(x)) \cdot (f^{-1}(x))' = 1$$

\textbf{Leibniz-Notation:}
$$\left.\frac{df^{-1}}{dx}\right|_{x=x_0} = \left(\left.\frac{df}{dx}\right|_{x=f^{-1}(x_0)}\right)^{-1}$$
\end{mainbox}

\begin{subbox}{Beispiel: Ableitung von $\ln(x)$}
	Mit $f(x) = \ln(x)$ und $g(x) = f^{-1}(x) = e^x$, $g'(x) = e^x$:
	\begin{align*}
		(\ln(x))' &= \frac{1}{g'(g^{-1}(x))} = \frac{1}{e^{\ln(x)}} = \frac{1}{x}
	\end{align*}
\end{subbox}

\begin{subbox}{Beispiel: Ableitung von $\sqrt{x}$}
	Mit $f(x) = \sqrt{x}$ und $g(x) = f^{-1}(x) = x^2$, $g'(x) = 2x$:
	\begin{align*}
		(\sqrt{x})' &= \frac{1}{g'(g^{-1}(x))} = \frac{1}{2\sqrt{x}}
	\end{align*}
	\textbf{Alternativ mit Potenzregel:}
	$(\sqrt{x})' = (x^{1/2})' = \frac{1}{2} x^{-1/2} = \frac{1}{2\sqrt{x}}$
\end{subbox}

\begin{mainbox}{Ableitungsregeln -- Zusammenfassung}
\renewcommand*{\arraystretch}{1.8}
\begin{center}
\begin{tabularx}{\linewidth}{l>{\centering\arraybackslash}X>{\centering\arraybackslash}X}
\toprule
\textbf{Regel} & \textbf{Newton} & \textbf{Leibniz} \\
\midrule
Konst.\ Summand & $(f + c)' = f'$ & $\frac{d(f+c)}{dx} = \frac{df}{dx}$ \\
Konst.\ Faktor & $(c \cdot f)' = c \cdot f'$ & $\frac{d(c \cdot f)}{dx} = c \cdot \frac{df}{dx}$ \\
Summenregel & $(f \pm g)' = f' \pm g'$ & $\frac{d(f+g)}{dx} = \frac{df}{dx} + \frac{dg}{dx}$ \\
Produktregel & $(f \cdot g)' = f' g + f g'$ & $\frac{d(fg)}{dx} = \frac{df}{dx} g + f \frac{dg}{dx}$ \\
Quotientenregel & $\left(\frac{f}{g}\right)' = \frac{f'g - fg'}{g^2}$ & $\frac{d}{dx}\left(\frac{f}{g}\right) = \frac{1}{g^2}\left(\frac{df}{dx} g - f \frac{dg}{dx}\right)$ \\
Kettenregel & $(f(g))' = f'(g) \cdot g'$ & $\frac{d(f \circ g)}{dx} = \frac{df}{dg} \cdot \frac{dg}{dx}$ \\
Umkehrfunktion & $(f^{-1})' = \frac{1}{f'(f^{-1})}$ & $\frac{df^{-1}}{dx} = \left(\frac{df}{dx}\big|_{f^{-1}}\right)^{-1}$ \\
\bottomrule
\end{tabularx}
\end{center}
\end{mainbox}

\subsection{Tangentengleichung}

\begin{mainbox}{Gleichung der Tangente}
	Tangente an $f(x)$ bei $x_0$ mit Steigung $f'(x_0)$:
	$$t(x) = f'(x_0) \cdot (x - x_0) + f(x_0)$$
	\textbf{Interpretation:} Die Tangente berührt den Graphen von $f$ bei $(x_0, f(x_0))$ und hat dort die gleiche Steigung wie $f$.
\end{mainbox}

\subsection{Begriffe und Interpretationen}

\begin{subbox}{Wichtige Begriffe}
	\begin{itemize}
		\item \textbf{Sekante:} Gerade durch zwei Punkte auf dem Funktionsgraphen
		\item \textbf{Tangente:} Gerade, die den Funktionsgraphen in einem Punkt berührt
		\item \textbf{Differenzenquotient:} Steigung der Sekante (durchschnittliche Änderungsrate)
		\item \textbf{Differentialquotient:} Grenzwert des Differenzenquotienten (momentane Änderungsrate)
		\item \textbf{Ableitung:} Tangentensteigung = Differentialquotient = $f'(x_0)$
		\item \textbf{Ableitungsfunktion:} Funktion, die jeder Stelle die Ableitung zuordnet
	\end{itemize}
\end{subbox}

\section{Folgen}
\subsection{Konvergenz}
Eine Folge $(a_n)_{n\in \mathbb{N}}$ konvergiert gegen L \\ $\iff \lim_{n \to \infty} a_n = L $
\\ $\iff \forall \varepsilon > 0 \ \exists N_\varepsilon \ \forall n \ge N_\varepsilon : \ | a_n - L | < \varepsilon$
\\ $\iff \forall \varepsilon > 0 $ ist die Menge $\{n \in \N \mid a_n \notin ]l-\varepsilon,l+\varepsilon[\}$ endlich.
\\ Es gibt für eine konvergente Folge $(a_n)_{n \geq 1}$ höchstens ein $L \in \R$, dass diese Bedingungen erfüllt. (Eindeutigkeit)


Wir dürfen (o.B.d.A.) annehmen, dass $\varepsilon $ durch eine Konstante $C \in \R$ beschränkt ist.
Es gilt ausserdem:
\begin{itemize}
 \item konvergent $\implies$ beschränkt, aber nicht umgekehrt
 \item $(a_n)$ konvergent $\iff (a_n)$ beschränkt \textbf{und} $\lim \inf a_n = \lim \sup a_n$
\end{itemize}

Sei $(a_n)_{n \geq 1}$, $(b_n)_{n \geq 1}$ konvergente Folgen. mit \\$\limn a_n = a, \limn b_n = b$ dann gilt:
\begin{itemize}
  \item $(a_n \pm b_n)_{n \geq 1}$ konvergent, $\limn (a_n \pm b_n) = a \pm b$
  \item $(a_n \cdot b_n)_{n \geq 1}$ konvergent, $\limn (a_n \cdot b_n) = a \cdot b$
  \item $(a_n \div  b_n)_{n \geq 1}$ konvergent, $\limn (a_n \div b_n) = a \div b$ ($b_n \neq 0, \forall n \in \N, b \neq 0$)
  \item Falls $\exists K \geq 1$ mit $a_n \leq b_n ~ \forall n \geq K$ folgt $a \leq b$
\end{itemize}

\begin{subbox}{Limes superior \& inferior}
  \vspace{-4pt}
$\limn \inf x_n = \limn \left( \inf_{m \ge n} x_m \right)$ \\
$\limn \sup x_n = \limn \left( \sup_{m \ge n} x_m \right)$
  \vspace{-12pt}
\end{subbox}

\begin{mainbox}{Einschliessungskriterium \\ (Sandwich-Theorem)}
  \vspace{-4pt}
Wenn $\limn a_n = \alpha, \ \limn b_n = \alpha$ und $a_n \le c_n \le b_n, \forall n \ge k$, dann $\limn c_n = \alpha$.
  \vspace{-4pt}
\end{mainbox}

\begin{mainbox}{Weierstrass}
Wenn $a_n$ monoton wachsend und nach oben beschränkt ist, dann konvergiert $a_n$ mit Grenzwert $\limn a_n = \sup \{a_n : \ n \ge 1\}$.

Wenn $a_n$ monoton fallend und nach unten beschränkt ist, dann konvergiert $a_n$ mit Grenzwert $\limn a_n = \inf \{a_n : \ n \ge 1\}$.
\end{mainbox}

\textbf{Tipp:} umschreiben $f(x)^{g(x)} = e^{g(x) \cdot \ln \left( f(x) \right)}$ und Grenzwert im Exponent berechnen (l'Hopital).
\begin{subbox}{	Berechne $\lim_{x \to 0^+} x^{\sin(x)}$.}
  \vspace{-6pt}
		$\lim_{x \to 0^+} x^{\sin(x)} = \lim_{x \to 0^+} e^{ \ln \left( x^{\sin(x)} \right)}
    = \lim_{x \to 0^+} e^{\sin (x) \ln (x)}$
	\\$e^x$ ist stetig $\Rightarrow$ Limes in Exponent.
	$
		\lim_{x \to 0^+} \sin(x) \ln(x) = \lim_{x \to 0^+} \frac{\ln(x)}{\frac{1}{\sin(x)}}
		\\\overset{\text{H}}{=} \lim_{x \to 0^+} \frac{\frac{1}{x}}{\frac{-\cos (x)}{\sin^2 (x)}} = \lim_{x \to0^+} \frac{- \sin^2 (x)}{x \cos (x)}\\
						= \lim_{x \to 0^+} \frac{- \sin x}{x} \cdot \frac{\overset{\to 0}{\sin x}}{\underset{\to 1}{\cos x}} = 0
	$
	\\Daraus folgt: $\lim_{x \to 0^+} e^{\sin (x) \ln (x)} = e^0 = 1$
  \vspace{-4pt}
\end{subbox}

\begin{mainbox}{Bernoulli Ungleichung}
  \vspace{-4pt}
  $(1 + x)^n \geq 1 + n \cdot x \qquad \forall n \in \N, x > -1$
  \vspace{-4pt}
\end{mainbox}

\begin{mainbox}{Cauchy-Kriterium}
  \vspace{-4pt}
Die Folge $a_n$ ist genau dann konvergent, falls \\ $\forall \varepsilon > 0, \ \exists N \ge 1$ so dass $\forall n,m \ge N: \\ | a_n - a_m | < \varepsilon$.
  \vspace{-4pt}
\end{mainbox}

\subsubsection{Teilfolge}
Eine Teilfolge von $a_n$ ist eine Folge $b_n$ wobei $b_n = a_{l(n)}$.
\\ $l$ ist eine Funktion mit $l(n) < l(n+1) \quad \forall n \ge 1$ (z.B. $l = 2n$ für jedes gerade Folgenglied).

\begin{subbox}{Bolzano-Weierstrass}
Jede beschränkte Folge besitzt eine konvergente Teilfolge.
\end{subbox}
 \textbf{Wenn 2 Teilfolgen nicht den gleichen Grenzwert haben, impliziert das Divergenz.}

\subsection{Strategie - Konvergenz von Folgen}
\begin{enumerate}
 \item Bei Brüchen: Grösste Potenz von $n$ kürzen. Alle Brüche der Form $\frac{a}{n^a}$ streichen, da diese nach 0 gehen.
 \item Bei Wurzeln in Summe im Nenner: Multiplizieren des Nenners und Zählers mit der Differenz der Summe im Nenner. (z.B. $(a+b)$ mit $(a-b)$ multiplizieren)
 \item Bei rekursiven Folgen: Anwendung von Weierstrass zur monotonen Konvergenz
 \item Einschliessungskriterium (Sandwich-Theorem) anwenden.
 \item Mit bekannter Folge vergleichen.
 \item Grenzwert durch einfaches Umformen ermitteln.
 \item Limit per Definition der Konvergenz zeigen.
 \item Anwendung des Cauchy-Kriteriums.
 \item Suchen eines konvergenten Majoranten.
 \item Weinen und die Aufgabe überspringen;)
\end{enumerate}

\subsection{Strategie - Divergenz von Folgen}
\begin{enumerate}
 \item Suchen einer divergenten Vergleichsfolge.
 \item Alternierende Folgen: Zeige, dass Teilfolgen nicht gleich werden, also $\limn a_{p_1(n)} \ne \limn a_{p_2(n)}$ (mit z.B. gerade/ungerade als Teilfolgen).
\end{enumerate}

\subsection{Tricks für Grenzwerte}
\subsubsection{Binome}
$$\lim_{x\to\infty} (\sqrt{x + 4} - \sqrt{x - 2}) = \lim_{x\to\infty} \frac{(x+4)-(x-2)}{\sqrt{x+4}+\sqrt{x-2}}$$

\subsubsection{Substitution}
$$\lim_{x\to\infty} x^2 \left(1-\cos\left(\frac{1}{x}\right)\right)$$
Substituiere nun $u = \frac{1}{x}$:
$$\lim_{u \to 0} \frac{1 - \cos(u)}{u^2} = \lim_{u \to 0} \frac{\sin(u)}{2u} = \lim_{u\to 0} \frac{\cos(u)}{2} = \frac{1}{2}$$

\subsection{Grenzwerte von Funktionen}
\begin{subbox}{Häufungspunkt}
  \vspace{-4pt}
 $x_0 \in \R$ ist ein Häufungspunkt der Menge D falls $\forall \delta > 0: (]x_0 - \delta, x_0 + \delta[ \backslash \{x_0\}) \cap D \ne \varnothing$.
 \vspace{-8pt}
\end{subbox}

\begin{mainbox}{Grenzwert - Funktionen}
  \vspace{-4pt}
 Wenn $f: D \to \R, x_0 \in \R$ ein Häufungspunkt von $D$ ist, dann ist $A \in \R$ der Grenzwert von $f(x)$ für $x \to x_0$ ($\lim_{x\to x_0} f(x) = A$), falls $\forall \varepsilon > 0 \ \exists \delta > 0$, so dass $\forall x \in D \cap (]x_0 - \delta, x_0 + \delta[ \backslash \{x_0\}): |f(x) - A| < \varepsilon$.
 \vspace{-4pt}
\end{mainbox}
Eine Funktion $f$ ist in $x_0$ stetig $\iff$ $\lim_{x \to x_0} f(x)$ existiert und gleich $f(x_0)$.
\\ \\
Die Stetigkeit einer abschnittweise definierten Funktion hängt auch vom Verhalten(links-/rechtsseitiger Grenzwert) an den Grenzen von Abschnittsintervallen ab.

\begin{subbox}{Satz von L'Hôpital}
  \vspace{-4pt}
  Seien $f,g$ stetig und differenzierbar auf $]a,b[$. Wenn $\lim_{x\to c} f(x) = \lim_{x \to c} g(x) = 0$ oder $\pm \infty$ und $g'(x) \ne 0 \ \forall x \in I \backslash \{c\}$, dann gilt $$\lim_{x\to c} \frac{f(x)}{g(x)} = \lim_{x\to c}\frac{f'(x)}{g'(x)}$$
  \vspace{-12pt}
\end{subbox}

Grenzwerte der Form $\infty^0$ und $1^\infty$ können meist mit $f(x)^{g(x)} = e^{g(x)\cdot \ln(f(x))}$ und dann Bernoulli (nur Exponenten betrachten da $e$ stetig) anwenden oder vereinfachen berechnet werden.

\begin{mainbox}{Wichtige Werte}
\begin{center}
 \begin{tabular}{c|cccccc}
  deg & 0° & 30° & 45° & 60° & 90° & 180° \\
  \midrule
  rad & 0 & $\frac{\pi}{6}$ & $\frac{\pi}{4}$ & $\frac{\pi}{3}$ & $\frac{\pi}{2}$ & $\pi$ \\
  cos & 1 & $\frac{\sqrt{3}}{2}$ & $\frac{\sqrt{2}}{2}$ & $\frac{1}{2}$ & 0 & -1 \\
  sin & 0 & $\frac{1}{2}$ & $\frac{\sqrt{2}}{2}$ & $\frac{\sqrt{3}}{2}$ & 1 & 0 \\
  tan & 0 & $\frac{1}{\sqrt{3}}$ & 1 & $\sqrt{3}$ & $+\infty$ & 0 \\
 \end{tabular}
\end{center}
\end{mainbox}

\begin{center}
% \includegraphics[width=\linewidth]{degrees_circle.pdf}
% Note: degrees_circle.pdf image file is required here
\end{center}

\section{Tabellen}
\subsection{Grenzwerte}
\arraycolsep=1pt\def\arraystretch{1.5}
\begin{center}
  \begin{tabularx}{\linewidth}{XX}
    \toprule
    $\limxi \frac{1}{x} = 0$ & $\limxi 1 + \frac{1}{x} = 1$ \\
    $\limxi e^x = \infty$ & $\limxn e^x = 0$ \\
    $\limxi e^{-x} = 0$ & $\limxn e^{-x} = \infty$ \\
    $\limxi \frac{e^x}{x^m} = \infty$ & $\limxn xe^x = 0$ \\
    $\limxi \ln(x) = \infty$ & $\limxo \ln(x) = -\infty$ \\
    $\limxi (1+x)^{\frac{1}{x}} = 1$ & $\limxo (1+x)^{\frac{1}{x}} = e$ \\
    $\limxi (1+\frac{1}{x})^b = 1$ & $\limxi n^{\frac{1}{n}} = 1$ \\
    $\lim_{x\to\pm\infty} (1 + \frac{1}{x})^x = e$ & $\limxi (1-\frac{1}{x})^x = \frac{1}{e}$ \\
    $\lim_{x\to\pm\infty} (1 + \frac{k}{x})^{mx} = e^{km}$ & $\limxi (\frac{x}{x+k})^x = e^{-k}$ \\
    $\limxo \frac{a^x -1}{x} = \ln(a), \forall a > 0$ &
    $\limxi x^a q^x = 0, \newline \forall 0 \le q < 1$ \\
  \end{tabularx}
  \begin{tabularx}{\linewidth}{XX}
    $\limxo \frac{\sin x}{x} = 1$ & $\limxo \frac{\sin kx}{x} = k$\\
    $\limxo \frac{1}{\cos x} = 1$ & $\limxo \frac{\cos x -1}{x} = 0$ \\
    $\limxo \frac{\log 1 - x}{x} = -1$ & $\limxo x \log x = 0$\\
    $\limxo \frac{1 - \cos x}{x^2} = \frac{1}{2}$ & $\limxo \frac{e^x-1}{x} = 1$ \\
    $\limxo \frac{x}{\arctan x} = 1$ & $\limxi \arctan x = \frac{\pi}{2}$ \\
    $\limxo \frac{e^{ax}-1}{x} = a$ & $\limxo \frac{\ln(x+1)}{x} = 1$ \\
    $\lim_{x\to 1} \frac{\ln(x)}{x-1} = 1$ & $\limxi \frac{\log(x)}{x^a} = 0$ \\
    $\limxi \sqrt[x]{x} = 1$ & $\limxi \frac{2x}{2^x} = 0$ \\
    \bottomrule
  \end{tabularx}
\end{center}
\subsection{Ableitungen}
\begin{center}
  % the c>{\centering\arraybackslash}X is a workaround to have a column fill up all space and still be centered
  \begin{tabularx}{\linewidth}{c>{\centering\arraybackslash}Xc}
  \toprule
  $\mathbf{F(x)}$ & $\mathbf{f(x)}$ & $\mathbf{f'(x)}$ \\
  \midrule
  $\frac{x^{-a+1}}{-a+1}$ & $\frac{1}{x^a}$ & $\frac{a}{x^{a+1}}$ \\
  $\frac{x^{a+1}}{a+1}$ & $x^a \ (a \ne 1)$ & $a \cdot x^{a-1}$ \\
  $\frac{1}{k \ln(a)}a^{kx}$ & $a^{kx}$ & $ka^{kx} \ln(a)$ \\
  $\ln |x|$ & $\frac{1}{x}$ & $-\frac{1}{x^2}$ \\
  $\frac{2}{3}x^{3/2}$ & $\sqrt{x}$ & $\frac{1}{2\sqrt{x}}$\\
  $-\cos(x)$ & $\sin(x)$ & $\cos(x)$ \\
  $\sin(x)$ & $\cos(x)$ & $-\sin(x)$ \\
  $\frac{1}{2}(x-\frac{1}{2}\sin(2x))$ & $\sin^2(x)$ & $2 \sin(x)\cos(x)$ \\
  $\frac{1}{2}(x + \frac{1}{2}\sin(2x))$ & $\cos^2(x)$ & $-2\sin(x)\cos(x)$ \\
  \multirow{2}*{$-\ln|\cos(x)|$} & \multirow{2}*{$\tan(x)$} & $\frac{1}{\cos^2(x)}$  \\
  & & $1 + \tan^2(x)$ \\
  $\cosh(x)$ & $\sinh(x)$ & $\cosh(x)$ \\
  $\log(\cosh(x))$ & $\tanh(x)$ & $\frac{1}{\cosh^2(x)}$ \\
  $\ln | \sin(x)|$ & $\cot(x)$ & $-\frac{1}{\sin^2(x)}$ \\
  $\frac{1}{c} \cdot e^{cx}$ & $e^{cx}$ & $c \cdot e^{cx}$ \\
  $x(\ln |x| - 1)$ & $\ln |x|$ & $\frac{1}{x}$ \\
  $\frac{1}{2}(\ln(x))^2$ & $\frac{\ln(x)}{x}$ & $\frac{1 - \ln(x)}{x^2}$ \\
  $\frac{x}{\ln(a)} (\ln|x| -1)$ & $\log_a |x|$ & $\frac{1}{\ln(a)x}$ \\
  \bottomrule
  \end{tabularx}
\end{center}
\subsection{Weitere Ableitungen}
\begin{center}
  \begin{tabularx}{\linewidth}{>{\centering\arraybackslash}X>{\centering\arraybackslash}X}
  \toprule
  $\mathbf{F(x)}$ & $\mathbf{f(x)}$ \\
  \midrule
  $\arcsin(x)$ & $\frac{1}{\sqrt{1 - x^2}}$ \\
  $\arccos(x)$ & $\frac{-1}{\sqrt{1 - x^2}}$ \\
  $\arctan(x)$ & $\frac{1}{1 + x^2}$ \\
  $x^x \ (x > 0)$ & $x^x \cdot (1 + \ln x)$ \\
  \bottomrule
  \end{tabularx}
\end{center}
\subsection{Integrale}
\begin{center}
 \begin{tabularx}{\linewidth}{>{\centering\arraybackslash}X>{\centering\arraybackslash}X}
  \toprule
  $\mathbf{f(x)}$ & $\mathbf{F(x)}$ \\
  \midrule
  $\int f'(x) f(x) \dx$ & $\frac{1}{2}(f(x))^2$ \\
  $\int \frac{f'(x)}{f(x)} \dx$ & $\ln|f(x)|$ \\
  $\int_{-\infty}^\infty e^{-x^2} \dx$ & $\sqrt{\pi}$ \\
  $\int (ax+b)^n \dx$ & $\frac{1}{a(n+1)}(ax+b)^{n+1}$ \\
  $\int x(ax+b)^n \dx$ & $\frac{(ax+b)^{n+2}}{(n+2)a^2} - \frac{b(ax+b)^{n+1}}{(n+1)a^2}$ \\
  $\int (ax^p+b)^n x^{p-1} \dx$ & $\frac{(ax^p+b)^{n+1}}{ap(n+1)}$ \\
  $\int (ax^p + b)^{-1} x^{p-1} \dx$ & $\frac{1}{ap} \ln |ax^p + b|$ \\
  $\int \frac{ax+b}{cx+d} \dx$ & $\frac{ax}{c} - \frac{ad-bc}{c^2} \ln |cx +d|$ \\
  $\int \frac{1}{x^2+a^2} \dx$ & $\frac{1}{a} \arctan \frac{x}{a}$ \\
  $\int \frac{1}{x^2 - a^2} \dx$ & $\frac{1}{2a} \ln\left| \frac{x-a}{x+a} \right|$ \\
  $\int \sqrt{a^2+x^2} \dx $ & $\frac{x}{2}f(x) + \frac{a^2}{2}\ln(x+f(x))$ \\
  \bottomrule
 \end{tabularx}
\end{center}

\begin{mainbox}{Integralformel}
  \vspace{-12pt}
  $$\int_a^b \frac{f(x)}{f(a+b-x)+f(x)}dx = \frac{b-a}{2}$$
  Beweis:
  $\int_a^b \frac{f(x)}{f(a+b-x)+f(x)}dx = I$
  \\Substitution: $u = a + b -x, du = -dx$
  $$I = \int_b^a \frac{f(a+b-u)}{f(u)+f(a+b-u)}(-du) $$
  $$= \int_a^b \frac{f(a+b-u)}{f(u)+f(a+b-u)}du$$
  $$2 \cdot I = \int_a^b \frac{f(x)+f(a+b-x)}{f(a+b-x)+f(x)}dx$$
  $$2I = \int_a^b1dx = b-a \implies I = \frac{b-a}{2}$$
\end{mainbox}

\subsection{Diverses}
\begin{itemize}
    \item $e=2,718$
    \item $ln(0)=1$
\end{itemize}
\section{Aufgaben}
\subsection{Multiple-Choice}%
\subsubsection{Generelle Tipps}

\begin{itemize}
	\item \textbf{Richtig lesen!!!}
	\item (Gegen-)Beispiele suchen.
	\item Monotonieverhalten bei verketteten/miteinander verechneten Funktionen bleibt nicht umbedingt erhalten.
    \begin{itemize}
      \item Sei $g: X \to Y$ streng monoton wachsend
      \begin{itemize}
        \item $f: Y \to Z$ (streng) monoton wachsend $\implies$ $f(g(x))$ (streng) monoton wachsend.
        \item $f: Y \to Z$ (streng) monoton fallend $\implies$ $f(g(x))$ (streng) monoton fallend.
      \end{itemize}
      \item Sei $g: X \to Y$ monoton wachsend
      \begin{itemize}
        \item $f: Y \to Z$ (streng) monoton wachsend $\implies$ $f(g(x))$ monoton wachsend.
        \item $f: Y \to Z$ (streng) monoton fallend $\implies$ $f(g(x))$ monoton fallend.
      \end{itemize}
      \item Die Inverse einer monoton wachsenden Funktion ist monoton wachsend.(Beweis mit Umkehrregel)
      \item Die Multiplikation oder Division zweier monotonen Funktionen können monoton wachsend/fallend bleiben oder ihre Monotonie verlieren.
      \item Monoton wachsende Funktion konvex(konkav) $\implies$ Inverse konkav(konvex).
      \item Monoton fallende Funktion konvex(konkav) $\implies$ Inverse konvex(konkav).
    \end{itemize}
	\item Verkettete Funktionen:
		\begin{itemize}
			\item Falls äussere beschränkt $\implies$ Verkettung beschränkt.
			\item Verkettung stetiger Funktionen ist stetig.
			\item Verkettung (streng) konvexer Funktionen, muss nicht konvex sein.
		\end{itemize}
	\item Schauen ob Funktionsfolge $f_n$ \underline{gleichmässig} und nicht nur punktweise gegen $f$ konvergiert.
		$ \implies$ Gewisse Eigenschaften (Stetigkeit) nur dann gültig.
	\item $f''(x_0) = 0 \centernot \implies x_0$ ist ein Sattelpunkt von $f$.
	\item $f$ stetig $\centernot \implies f'$ stetig.
	\item Sei $f : \R \to \R$ eine ungerade Funktion. $\implies$ $f^{(i)} (0) = 0$ für $i$ gerade.
	\item $X,Y,Z \subset \R$, $f: X \to Y$ und $g: Y \to Z$, sodass $g \circ f : X \to Z$ bijektiv.
		$\implies$ $f$ injektiv, $g$ surjektiv
	\item Sei $a,b \in \R$ mit $a < b$ mit $f: ]a,b[ \to \R$. Falls $f^2$ und $f^3$ in $]a,b[$
		differenzierbar und $f(x) \neq 0~\forall x \in ]a,b[$ $ \implies$ $f$ in $]a,b[$ differenzierbar. (Da wir $f^3$ durch $f^2$ teilen).
	\item Im Allg. muss eine Grenzfunktion $f$ nicht differenzierbar sein, und wenn sie es ist muss $f'$ nicht gleich $\limn f_n'$ sein.

	\item $f: (0,1] \to \R$ monoton $\centernot \implies f$ beschränkt.
	\\ Der Logarithmus $\ln: (0,1] \to \R$ ist monoton, aber nicht beschränkt.
  \item $f: [a, b] \to \R$ für beliebige $a, b \in \R, a < b$ ist beschränkt.
  \item $f(x) = |x| \cdot \text{sign}(x)$ ist stetig, da $|x| \cdot \text{sign}(x) = x$ und $g(x)=x$ stetig.
\end{itemize}

\subsubsection{Gegenbeispiele}

    \paragraph{Reelle Zahlen}
    \begin{itemize}
        \item Die rationalen Zahlen $\mathbb{Q}$ sind kein vollständiger Körper.
        \item Die komplexen Zahlen $\mathbb{C}$ sind kein angeordneter Körper.
        \item Die reellen Zahlen $\mathbb{R}$ sind ein vollständig angeordneter Körper.
        \item Eine Menge, die kein Maximum besitzt: $A = \lbrace{x \in \mathbb{Q}  \vert x \le \sqrt{5}}\rbrace$
    \end{itemize}
    \paragraph{Folgen}
    \begin{itemize}
        \item Seien $a_{n} = \frac{1}{n^2}$ und $b_{n} = \frac{1}{n}$. Dann gilt $a_{n} \le b_{n}, \forall n \in \mathbb{N}$, aber für die Grenzwerte $lim_{n \to +\infty} a_{n} = a$ und $lim_{n \to +\infty} b_{n} = b$ gilt $a \le b$, da $a=b=0$.
        \item Die Folge $(-1)^n$ ist beschränkt, aber nicht konvergent, insbesondere weil $\liminf_{n\to \infty}{b_{n}} = -1$ ungleich $\limsup_{n \to \infty}{b_{n}} = 1$ ist.
        \item Sei $b_{n} = (-1)^n$. Dann ist $b_{n}^2$ konvergent, $b_{n}$ aber nicht.
        \item Die Folge $1,2,2,2+\frac{1}{2},3,3+\frac{1}{3},3+\frac{2}{3},4,4,4+\frac{1}{4},4+\frac{2}{4},4+\frac{3}{4},5...$ erfüllt das Cauchy-Kriterium für Folgen nicht, da nur $\vert a_{n}-a_{n+1}\vert < \epsilon$ für $n \ge N$.
    \end{itemize}
    \paragraph{Reihen}
    \begin{itemize}
        \item Eine Reihe die konvergiert aber nicht absolut: $\displaystyle\sum_{n=1}^{\infty} (-1)^n \frac{1}{n}$ (alternierende harmonische Reihe)
        \item Die harmonische Reihe $\displaystyle\sum_{n=1}^{\infty} \frac{1}{n}$ divergiert, aber $\displaystyle\sum_{n=1}^{\infty} \frac{1}{n^p}$ konvergiert für $p>1$.
    \end{itemize}
    \paragraph{Stetigkeit und Funktionenfolgen}
    \begin{itemize}
        \item Eine Stufenfunktion, wie bspw.
        \[f(x)= \begin{cases}
                1 &\text{für } x \ge 0 \\
                0 &\text{für } x < 0
                \end{cases}\]
        ist unstetig im Punkt $0$.
        \item Die Funktion
        \[f(x)= \begin{cases}
                sin(\frac{1}{x}) &\text{für } x \in \mathbb(R) \\
                0 &\text{für } x=0
                \end{cases}\]
        ist unstetig in $0$.
        \item Die Funktionenfolge $f_{n} = x^n$ auf $[0,1]$ konvergiert punktweise gegen
        \[f(x)= \begin{cases}
                0 &\text{für } x \in ]0,1[ \\
                1 &\text{für } x=1
                \end{cases}\]
        aber die Konvergenz ist nicht gleichmässig gegen $1$ bei $x=1$.
        \item Die Funktion $f(x)x^2$ ist stetig, aber nicht gleichmässig stetig.
    \end{itemize}
    \paragraph{Differentialrechnung}
    \begin{itemize}
        \item Die Funktion $f(x)=\vert x-a \vert$ für ein $a \in \mathbb(R)$ ist nicht differenzierbar im Punkt $x=a$.
        \item Die Weierstrass-Funktion ist überall stetig, aber nirgends differenzierbar auf $\mathbb{R}$.
        \item Sei $f(x)=x^3$. Dann gilt $f'(0) = 3*0^2$, aber $f$ hat in $0$ kein lokales Extremum.
        \item Sei $f(x)=x^3$, also ist f streng monoton fallend. Aber es gilt $f'(x) = 3x^2 \ge 0$ für alle $x \in \mathbb{R}$ und nicht $f'(x) > 0$.
        \item Unstetige Ableitung: Sei
        \[f(x)= \begin{cases}
                x^2 sin(\frac{1}{x}) &\text{für } x \ne 0 \\
                0 &\text{für } x = 0
                \end{cases}\]
        Dann hat $f$ in $0$ eine Ableitung die unstetig ist.
        \item Sei $f(x) = \frac{2}{x} + cos(\frac{1}{x})$ und $g(x) = \frac{1}{x}$. Dann existiert $\lim_{x \to 0}{\frac{f(x)}{g(x)}}$, aber $\lim_{x \to 0}{\frac{f'(x)}{g'(x)}}$ nicht.
    \end{itemize}
    \paragraph{Integralrechnung}
    \begin{itemize}
        \item Die Dirichletfunktion
        \[f(x)= \begin{cases}
                1 &\text{für } x \in \mathbb{Q} \\
                0 &\text{für } x \in \mathbb{R} \backslash \mathbb{Q}
                \end{cases}\]
        ist nicht stetig aber trotzdem Riemann-integrierbar.
        \item Die Dirac-$\delta$ Funktion ist nicht Riemann-integrierbar.
        \item Das Integral $\int_{1}^{\infty}{\frac{1}{x}dx}$ existiert nicht, aber für $a > 1$ exisitiert $\int_{1}^{\infty}{\frac{1}{x^a}dx}$.
    \end{itemize}
\subsection{Trigonometrie}
    Manipulation für $f(x) = \sin(\arctan (x))$ (funktioniert ähnlich für $\cos(\arctan (x))$):

    \begin{align*}
      f(x) &=\frac{\sin(\arctan(x))}{\sqrt{\sin^2(\arctan(x))+\cos^2(\arctan(x))}}\\
      &=\frac{\frac{\sin(\arctan(x))}{\cos(\arctan(x))}}{\sqrt{\frac{\sin^2(\arctan(x))}{\cos^2(\arctan(x))}+1}}\\
      &=\frac{\tan(\arctan(x))}{\sqrt{\tan^2(\arctan(x))+1}}\\
      &=\frac{x}{\sqrt{x^2+1}}
    \end{align*}

    Manipulation für $f(x) = \tan(\arcsin (x))$:
    \begin{align*}
      f(x)&=\tan(\arcsin (x))\\
      &=\frac{\sin(\arcsin(x))}{\cos(\arcsin(x))}\\
      &=\frac{\sin(\arcsin(x))}{\sqrt{1-\sin^2(\arcsin(x))}}\\
      &=\frac{x}{\sqrt{1-x^2}}.
      \end{align*}

\subsection{Konvergenzbeweise}
\subsubsection{Folgen}
    Beweise: $lim_{n \to \infty}\frac{n^2}{n^2 +1}=1$\\
    Wähle: $\epsilon = 10^{-5}$. Finde ein $N \in \mathbb{N}$, sodass $\forall n \ge N$ gilt:
    \[\lvert a_n - a\rvert=\lvert\frac{n^2}{n^2 +1}-1\rvert <\epsilon\]
    Forme um:
    \[\lvert \frac{n^2}{n^2 +1}-1\rvert = \lvert \frac{n^2 -n^2 -1}{n^2 +1}\rvert = \frac{1}{n^2 +1}<\epsilon\]
    \[\implies n^2 +1 > \frac{1}{\epsilon} \implies n>\sqrt{\frac{1}{\epsilon}-1}\]
    Die Bedingung $\lvert a_n -a\rvert <\epsilon$ impliziert somit $n>\sqrt{\frac{1}{\epsilon}-1}$. Damit eine natürliche Zahl entstehen kann, wählen wir $N \coloneqq [\sqrt{\frac{1}{\epsilon}-1}]$. Für alle $n\ge N$ gilt somit nach Konstruktion:
    \[\lvert \frac{n^2}{n^2 +1}-1\rvert <\epsilon\]
    Dies entspricht genau der anfänglichen Definition von $lim_{n \to \infty}\frac{n^2}{n^2 +1}=1$.
\vspace{0.1 cm}
\hrule
\vspace{0.2 cm}
    Man untersuche die nachstehenden Zahlenfolge. Konvergieret sie? Wenn ja: Welches ist ihr Grenzwert?
    \[lim_{n\to +\infty}\frac{2n^3+n^2+1}{5(n^3)!+2}\]
    Sei $\epsilon >0$. Es gilt:
    \[0<\frac{2}{5(n^3)!+2} <\frac{2n^3+n^2-1}{5(n^3)!+2}\]
    \[\le \frac{3n^3}{5(n^3)!+2}<\frac{3n^3}{3(n^3)!}=\frac{n^3}{(n^3)!}<\epsilon\]
    wenn $n>N$ ist. Die Folge konvergiert gegen 0.
\subsubsection{Reihen}
    Zeige, dass der Konvergenzradius der Reihe
    \[\sum_{n=0}^{+\infty} \frac{(n!)^2}{(2n)!}z^n\]
    gleich $4$ ist.
    Sei $x_n \coloneqq \frac{(n!)^2 z^n}{(2n)!}$. Die Reihe konvergiert trivialerweise, wenn $z=0$ ist. Nehmen wir an, dass $z\ne 0$ ist. Wir brechnen den folgenden Limes:
    \[lim_{n\to +\infty} \frac{\lvert x_{n+1}\rvert}{\lvert x_n\rvert}\]
    \[= lim_{n\to +\infty}\frac{((n+1)!)^2}{(2(n+1))!}\lvert z\rvert ^{n+1}\frac{(2n)!}{\lvert z\rvert ^n (n!)^2}\]
    \[lim_{n\to +\infty}\frac{(n!)^2(n+1)^2}{(2n)!(2n+1)(2n+2)}\frac{(2n)!}{(n!)^2}\lvert z\rvert\]
    \[=\lvert z\rvert lim_{n\to +\infty}\frac{(n+1)^2}{(2n+1)(2n+2)}\]
    \[\lvert z\rvert lim_{n\to +\infty}(\frac{1}{4}+\frac{1}{4(2n+1)}) =\frac{\lvert z\rvert}{4}\]
    Aus dem Quotientenkriterium können wir folgern, dass die Reihe absolut konvergent ist falls
    \[limsup_{n\to \infty}\frac{\lvert x_{n+1}\rvert}{\lvert x_n\rvert}\]
    \[=lim_{n\to +\infty}\frac{\lvert x_{n+1}\rvert}{x_n}=\frac{\lvert z\rvert}{4}<1\]
    wobei $\lvert z\rvert <4$. Die Reihe ist divergent, falls
    \[limsup_{n\to \infty}\frac{\lvert x_{n+1}\rvert}{\lvert x_n\rvert}\]
    \[=lim_{n\to +\infty}\frac{\lvert x_{n+1}\rvert}{x_n}=\frac{\lvert z\rvert}{4}>1\]
    wobei $\lvert z\rvert >4$. Somit haben wir gezeigt, dass der Konvergenzradius der Reihe gleich 4 ist.

\subsubsection{Funktionen}
    Beweise, dass folgende Gleichung gilt:
    \[lim_{n\to +\infty}(\sqrt{n}-\sqrt{n+1})=0\]
    Es gilt: $\sqrt{n}-\sqrt{n+1}=(\sqrt{n}-\sqrt{n+1})\frac{\sqrt{n}+\sqrt{n+1}}{\sqrt{n}+\sqrt{n+1}}$
    \[=\frac{n-(n+1)}{\sqrt{n}+\sqrt{n+1}}=-\frac{1}{\sqrt{n}+\sqrt{n+1}}\]
    \[\implies \lvert \sqrt{n}-\sqrt{n+1}\rvert =\frac{1}{\sqrt{n}+\sqrt{n+1}}\le \frac{1}{2\sqrt{n}}\]
    Sei $\epsilon >0$. Dann:
    \[\frac{1}{2\sqrt{n}}<\epsilon \iff \frac{1}{2\epsilon}<\sqrt{n}\iff \frac{1}{4{\epsilon}^2}<n\]
    Wir haben gezeigt, dass für alle $\epsilon >0$ die Ungleichung
    \[\lvert \sqrt{n}-\sqrt{n+1}\rvert <\epsilon\]
    gilt, falls $n>N \coloneqq (4{\epsilon}^2)^{-1}$

\subsection{Grenzwertberechnungen}
    \[lim_{x\to 6}\frac{sin(\pi x)}{x^2 -36}\]
    Benutze den Satz von L'Hôpital:
    \[sin(\pi x)' = \pi cos(\pi x); (x^2 -36)' = 2x\]
    \[\implies lim_{x \to 6}\frac{sin(\pi x)}{x^2 -36}=lim_{x \to 6}\pi \frac{cos(\pi x)}{2x}\]
    Ausserdem gilt: $cos(6\pi)=1$ und $2\cdot 6 = 12$
    \[\implies lim_{x\to 6}\pi \frac{cos(\pi x)}{2x} = \pi \frac{1}{12} = \frac{\pi}{12}\]
\vspace{0.1 cm}
\hrule
\vspace{0.2 cm}
    Bestimme: $lim_{n\to \infty} \sqrt{n^4+7n^2}-n^2$
    \[=lim_{n\to \infty} \frac{(\sqrt{n^4+7n^2}-n^2)\cdot (\sqrt{n^4+7n^2}-n^2)}{\sqrt{n^4+7n^2}-n^2}\]
    \[=lim_{n\to \infty} \frac{(n^4+7n^2)-n^4}{\sqrt{n^4(1+\frac{7}{n^2})}+n^2}\]
    \[=lim_{n\to \infty} \frac{7n^2}{n^2\sqrt{1+\frac{7}{n^2}}+n^2}\]
    \[=lim_{n\to \infty} \frac{7}{\underbrace{\sqrt{1+\frac{7}{n^2}}}_{\to 1}+1} = \frac{7}{2}\]
\vspace{0.1 cm}
\hrule
\vspace{0.2 cm}
    Bestimme: $lim_{x\to 6}\frac{sin(\pi x)}{x^2-36}$\\
    Da das Limit gegen 0 läuft, ist die Verwendung der Regel von L'Hôpital angezeigt:
    \[sin(\pi x)'=\pi cos(\pi x)\textit{; }(x^2-36)'=2x\]
    \[lim_{x\to 6}\frac{sin(\pi x)}{x^2 -36}=lim_{x\to 6}\frac{cos(\pi x)}{2x}\]
    Ausserdem wissen wir: $cos(6\pi)=1$ und $2\cdot 6=12$. Daraus folgt:
    \[lim_{x\to 6}\pi \frac{cos(\pi x)}{2x}=\pi \frac{1}{12}=\frac{\pi}{12}\]
\subsection{Differentialrechnungen}
\subsubsection{Produktregel}
    \[f(x)=xe^x\]
    \[f'(x)=e^x+xe^x=(x+1)e^x\]
\subsubsection{Quotientenregel}
    \[f(x)=(\frac{e^x}{x+1})\]
    \[f'(x)=\frac{(1+x)\frac{d}{dx}(e^x)-e^x\frac{d}{dx}(x+1)}{(x+1)^2}\]
    \[=\frac{(1+x)e^x -e^x}{(1+x)^2} = \frac{xe^x}{(1+x)^2}\]
\subsubsection{Kettenregel}
    \[f(x) = exp(x+exp(x))\]
    \[f'(x) = exp(x+exp(x))*(1+exp(x))\]
    \[= exp(x+exp(x))+exp(2x+exp(x))\]
\subsubsection{Umkehrregel}
    \[f(x)=log(x)\]
    \[y=e^x\implies x=log(y)\]
    \[y'=e^x\implies x'=\frac{1}{y'}=\frac{1}{e^x}=\frac{1}{e^{log(y)}}=\frac{1}{y}\]
    Somit ist $log(x)'=\frac{1}{x}$
\subsubsection{Taylor-Polynome}
    Bestimmen Sie mit Hilfe des linearen Taylorpolynoms um $t_0 = 8$ eine Näherung
    an $\sqrt[3]{7}$.
    \\Geben Sie eine Abschätzung für den Fehler dieser Näherung an.
    \\Für $f(t) = \sqrt[3]{t}$ ist $f'(t) = \frac{1}{3}t^{-\frac{2}{3}}$.
    \\$\implies$
    $$T_1(f, t, 8) = f(8) + f'(8)(t - 8) = 2 + \frac{1}{3}\left(\frac{1}{2}\right)^2 \cdot (t- 8)$$
    Näherung für $\sqrt[3]{7} = T_1(f, 7, 8) = 2 - \frac{1}{12} = \frac{23}{12}$
    \\Fehlerabschätzung:
    $$R_1(f, t, t_0) = \frac{f''(\xi)}{2!}(t - t_0)^2 = \frac{f''(\xi)}{2!}(7-8)^2 = \frac{f''(\xi)}{2}$$
    für ein $\xi \in [7, 8]$.
    Da $f''(t) =-\frac{2}{9} t^{-\frac{5}{3}}$ gilt:
    $$|\text{Fehler}| \leq \frac{1}{2} \sup_{\xi \in [7, 8]} |f''(\xi)| = \frac{1}{2}\cdot \frac{2}{9} \cdot \frac{1}{7^{5/3}} = \frac{1}{9 \cdot 7^{5/3}}.$$
\vspace{0.1 cm}
\hrule
\vspace{0.2 cm}
    Betrachte die Funktion $f: (0,+\infty)\to \mathbb{R}$ definiert durch:
    \[f(x)= \int_x^2 \frac{e^{t^2}}{t}dt + log(\frac{x}{2})\]
    (a) Berechne die Ableitung von $f$ und zeige, dass $f$ umkehrbar ist:\\
    \[f'(x) = \frac{1}{x}(1-e^{x^2})\]
    Für alle $x\in(0,+\infty)$ haben wir $e^{x^2} > 1$ und daher auch $\frac{1}{x}(1-e^{x^2}) < 0$ für alle $x \in (0,+\infty)$. Nun folgt aus dem Umkehrsatz, dass $f$ umkehrbar, und $f^{-1}$ differenzierbar ist.\\
    (b) Berechne $f(2)$ und bestimme das Taylor-Polynom 1. Ordnung von $f^{-1}$ an der Stelle $x_0 = 0$:\\
    \[f(2) = \int_2^2...+log(1) = 0\]
    Das Taylor-Polynom berechnet sich folgendermassen:
    \[f_{\textit{Taylor}}^{-1}(x) = f^{-1}(0)+(f^{-1})'(0)(x-0)\]
    Aus Teil (a) wissen wir, $f^{-1}(0) = 2$. Aufgrund von Korollar 4.1.12 wissen wir $(f^{-1})'(y_0) = \frac{1}{f'(x_0)}$.\\
    Daraus folgt: $(f^{-1})'(0) = \frac{1}{f'(2)} = \frac{2}{1-e^4}$.\\
    Nun erhalten wir: $f_{\textit{Taylor}}^{-1} = 2+\frac{2}{1-e^4}x$
\subsection{Integralrechnungen}
    Sonderfall: Ungerade Funktionen, bspw.
    \[\int_{-\frac{\pi}{2}}^{\frac{\pi}{2}}(sin(x))^7 cos(x)dx\]
    Diese Integrale sind stets $=0$.
\subsubsection{Partielle Integration}
    \[\int_{-\infty}^{0}xexp(3x)dx\]
    \[=[x\cdot \frac{1}{3}exp(3x)]_{-\infty}^{0} - \int_{-\infty}^{0}1 \cdot \frac{1}{3}exp(3x)dx\]
    \[=[x \cdot \frac{1}{3}exp(3x)]_{-\infty}^{0} - \frac{1}{3}[\frac{1}{3}exp(3x)]_{-\infty}^{0}\]
    \[=0-0-\frac{1}{3}(\frac{1}{3} \cdot 1-0) = -\frac{1}{9}\]
    aufgrund von $lim_{x\to \infty}[xexp(3x)] \to 0$
\subsubsection{Substitution}
    \begin{align*}
      \int x\arcsin(x)\,\mathrm dx\\
    \end{align*}
    Substitution mit $ x = \sin(u) $:
    \begin{align*}
      =\int \sin(u)\cos(u) u \, \mathrm du\\
      = \int \frac{1}{2} sin(2u) u \, \mathrm du\\
      = \frac{1}{2} \int \sin(2u) u \, \mathrm du\\
    \end{align*}
     Partielle Integration ( $u$ ableiten, $\sin(2u)$ integrieren):
     \begin{align*}
      \frac{1}{2} \int u \sin(2u)\, \mathrm du \\
      = \frac{1}{2} \left[ u \cos(2u)  \left(-\frac{1}{2}\right) - \int \left(-\frac{1}{2}\right) \cos(2u) \, \mathrm du \right] \\
      =-\frac{1}{4} u \cos(2u) + \frac{1}{4} \int \cos(2u) \, \mathrm du\\
      = -\frac{1}{4} u \cos(2u) + \frac{1}{8} \sin(2u) + C
    \end{align*}
    $u = \arcsin(x)$ einsetzen liefert:
    $$= -\frac{1}{4} \arcsin(x) \cos(2\arcsin(x)) + \frac{1}{8} \sin(2\arcsin(x)) + C$$
\vspace{0.1 cm}
\hrule
\vspace{0.2 cm}
    Berechne: $\int_0^x xcos(x^2)dx$\\
    Sei $u = x^2$, $du=2xdx \iff \frac{1}{2}du = xdx$\\
    \[\int_0^x xcos(x^2)dx = \int_0^x \frac{1}{2}cos(u)du\]
    \[=\frac{1}{2}[sin(u)]_{0}^{y}\]
    \[=\frac{1}{2}sin(y^2)\]
\vspace{0.1 cm}
\hrule
\vspace{0.2 cm}
    Integriere das folgende uneigentliche Integral:
    \[\int_1^{\infty} \frac{1}{x^3}\sqrt{\frac{x}{x+1}}dx\]
    Wir formen zuerst um\\
    und machen die folgende Abschätzung:
    \[\frac{1}{x^3}\sqrt{\frac{1}{x+1}}\le \frac{1}{x^3}\]
    Da $\frac{1}{x^3}$ bekanntlich konvergent ist und unser Integral monoton wachsend, können wir schliessen, das der Grenzwert für das Integral existiert, d.h. es konvergiert.\\
    Wir führen die folgende Substitution durch: $t=\sqrt{\frac{x}{x+1}}$ und erhalten:
    \[\int_1^{\infty}\sqrt{\frac{x}{x+1}}dx\]
    \[=2(-\frac{1}{3}(\frac{t}{t+1})^{-\frac{3}{2}}+\sqrt{\frac{t}{t+1}}+\frac{2^{\frac{3}{2}}}{3})-\sqrt{2}\]
    Für $t\to\infty$ ergibt das $\frac{3}{2}(2-\sqrt{2})$.
\subsubsection{Partialbruchzerlegung}
    \[\int_{-2}^{-1}\frac{5x-8}{x^3 -4x^2 +4x}dx\]
    \[=\int_{-2}^{-1}\frac{5x-8}{x(x-2)^2}dx\]
    Bestimme die Nullstellen:
    \[x_1 =0; x_2 =2\]
    \[\implies \frac{A}{x}+\frac{B}{x-2}+\frac{C}{(x-2)^2}\]
    Bestimme A durch Multiplikation mit $x(x-2)^2$ und Einsetzen von $x=0$
    \[\implies A=-2\]
    Bestimme C durch Multiplikation mit $x(x-2)^2$ und Einsetzen von $x=2$
    \[\implies C=1\]
    Bestimme B durch einsetzen von A und C und anschliessendem  aurechnen:
    \[\frac{B}{x-2}=\frac{5x-8}{x(x-2)^2}+\frac{2}{x}-\frac{1}{x-2}^2\]
    \[\frac{5x-8+2(x-2)^2 -x}{x(x-2)^2} = \frac{2}{x-2}\]
    \[\implies B=2\]
    Schreibe das durch die PBZ entstandene Integral auf:
    \[\int_{-2}^{-1}\frac{-2}{x}dx+\int_{-2}^{-1}\frac{2}{x-2}dx+\int_{-2}^{-1}\frac{1}{(x-2)^2}dx\]
    Berechne die einzelnen Integrale und setze die Grenzen ein:
    \[\lbrack(-2ln(\lvert x\rvert)+2ln(\lvert x-2\rvert)-\frac{1}{(x-2)^2})\rbrack_{-2}^{-1}\]
    \[2\cdot(ln(2)+ln(3)-ln(4))+\frac{1}{12}\]
    \[=2\cdot(ln(\frac{3}{2}))+\frac{1}{12}\]

\end{document}